\section{Verification}
\label{sec:verification}

% Owner: Christopher Rodolph (@Capheuss). Tracking: GitHub issue #26.
% Each paragraph maps to the GitHub issues noted in its comment; the
% milestone gates (M2-M4) are in issue #17 and project-docs/PLAN.md.

Section~\ref{sec:approach} describes how we measure each configuration; this
section describes how we establish that a configuration is worth measuring,
that is, that it still computes the right results and that it actually closes
the leak. STT only ever \emph{delays} instructions; it never changes what they
compute. This means every check below must determine that a protected core
produces exactly the architectural results of the baseline, and differs
from it only in timing. We verify in layers, from single modules up to the
full SoC, so that a bug is caught at the smallest level where it can appear.

% Issues #4 (A2), #5 (A3).
\paragraph{Module-level tests}
The rename-stage logic is tested in isolation before it is wired into the core.
Directed ChiselTest cases cover the situations where a load's
destination takes its own ROB index, the same-cycle bypass between
instructions renamed together, and restoring the taint state after a branch
mispredict or exception, all of which are corner cases that can occur in the taint table.
The Age Matrix is checked by equivalence rather than by hand-written cases. 
For randomly generated rename groups it must produce exactly the same YRoTs as the 
baseline bypass, and we reuse the random test generator released with 
$\mu$STT~\cite{Chen:uSTT:2025} for this.

% Issue #10 (C2): riscv-tests and Spike co-simulation in the regression script.
\paragraph{Functional regression}
Every configuration in Table~\ref{tab:configs} must pass the RISC-V ISA tests
and the riscv-tests benchmarks in Verilator, as the baseline does. Because a
delay mechanism can fail by stalling forever rather than by computing a wrong
value, a test that times out counts as a failure. This is crucial for catching
deadlocks such as a transmitter that waits on a taint that is never cleared. 
We also run Chipyard's co-simulation mode, which checks each committed instruction
against the Spike ISA simulator in lockstep, so a divergence is reported at
the first wrong instruction rather than at the end of a test. With every new
switch turned off, the generated Verilog must be identical to the baseline's,
which we check by diffing it; this confirms that our changes leave the
unmodified design untouched and baseline numbers stay comparable.

% Issues #9 (C1), #6 (B1), #7 (B2).
\paragraph{Security checks}
The Spectre proofs of concept are the end-to-end security tests which all print
whether it recovered the secret and enable the regression script to check the
result automatically. An attack must leak on the baseline, which shows the
test itself works, and must fail on every protected configuration. Since the
conservative stand-in is strictly stronger than STT, it gives a second reference point 
such that any attack that succeeds against it points to a path that bypasses the issue queue, 
rather than to a flaw in taint tracking. Because the attacks cover only a few gadgets, 
we also add a simulation assertion at every point where a transmitter can fire, 
including the load/store unit's retry, replay and forwarding paths. If a transmitter 
fires while its YRoT is not yet safe, simulation stops. This checks the STT rule on every 
instruction of every test, not only on the attack code.

% Issues #14 (D3), #15 (D4).
\paragraph{Synthesized design}
Finally, short targeted programs, such as the Spectre gadget, are run in Vivado
simulation of the synthesized design to confirm that the logic Vivado
produces behaves like the Chisel we verified in Verilator. Before that, the
testbenches released with $\mu$STT's standalone RTL are run in the same
simulator, which should reproduce the published results and validate our
simulation setup.

\paragraph{Limitations}
These checks are tests, not proofs. Passing them shows that the known Spectre
variants are blocked and that the assertion never fires on our workloads, but
not that no leak exists. Formal analysis of the RTL, as done in
SecureBOOM~\cite{Jauch:SecureBOOM:2023}, would give that stronger guarantee, but
that is outside the scope of this project.
